\documentclass[11pt]{article}
\usepackage[a4paper,margin=1in]{geometry}
\usepackage{amsmath,amsthm}
\usepackage{unicode-math}
\setmathfont{STIX Two Math}
\usepackage{microtype}
\usepackage[colorlinks=true,linkcolor=black,citecolor=black,urlcolor=black]{hyperref}
\theoremstyle{plain}
\newtheorem{theorem}{Theorem}[section]
\newtheorem{lemma}[theorem]{Lemma}
\newtheorem{proposition}[theorem]{Proposition}
\newtheorem{corollary}[theorem]{Corollary}
\theoremstyle{definition}
\newtheorem{definition}[theorem]{Definition}
\theoremstyle{remark}
\newtheorem*{remark}{Remark}
\setlength{\parskip}{0.35em}
\setlength{\parindent}{0pt}
\title{The forcing chain and the residue: a substitution, and the sign that survives everything it builds}
\author{Dritëro M.}
\date{2026}
\begin{document}
\maketitle
\begin{abstract}
The substitution \(\sigma \): a \(\to \) ab, b \(\to \) a generates the entire program: its eigenvalue is the golden ratio, its companion matrix is the golden monodromy, its mapping torus is the figure-eight, its trace field is \(\mathbb{Q}(\sqrt{-3})\), and iterating it walks the Markov spine. We follow \(\sigma \) up the tower it builds and ask what survives every level. The answer is a single sign. The orientation character — the norm \(N(\lambda _{m}) = -1 = det X_{m}\) of the metallic half-matrix, the determinant \(det(Par\cdot W(w)) = -\omega ^\{\#L-\#R\}\) of the quantized word, the \(\mathbb{Z}/2\) that separates the two-register breath det(Par@N) \(= (-1)^\{(N-1)/2\}\) — is one and the same residue, appearing in the arithmetic, the geometry, and the quantization as the one quantity that does not cancel. We prove the residue's laws: it is frozen through the \(\varphi -power\) degeneracy, it beats with the Pisano period \(\pi (3) = 8\) along the letter tower, and it is the fixed point of the object under its own half-turn — the \emph{held breath}, which we show is exactly the torsion: member m holds the breath of every divisor d \(\ge \) 3 of m. The program's method mirrors its subject. Cancel everything; what remains is real, and what remains is the sign.  ---
\end{abstract}
\section{Introduction}

A generating rule is a small thing that will not stop. \(\sigma \): a \(\to \) ab, b \(\to \) a is two arrows, and from them the whole object unfolds: the golden ratio as the growth rate, the matrix \(\left(\begin{smallmatrix}1 & 1\\1 & 0\end{smallmatrix}\right)\) as the linearization, the figure-eight as the mapping torus, \(\mathbb{Q}(\sqrt{-3})\) as the field, the Markov numbers as the traces. This paper is about what such a rule \emph{keeps}. Most of what a substitution produces cancels as you climb — traces grow, fields escalate, volumes accumulate, and the specific values wash out into class data. But underneath the growth one quantity is conserved at every level, and it is not a magnitude. It is a sign.

We call it the residue, and its recurring appearance is the paper's thesis: the norm of the scale \((-1)\), the determinant of the quantized word (\(-\omega \) to a power), the parity of the classical map \((\pm 1)\) are the same \(\mathbb{Z}/2\), seen through three different instruments. Section 2 lays out the forcing chain — how \(\sigma \) forces each next object with no freedom. Section 3 proves the norm identity and its freezing through the golden degeneracy. Section 4 gives the residue's determinant law and its Pisano rhythm. Section 5 is the two-register breath — the residue read simultaneously in the quantum and classical registers. Section 6 is the deepest: the \emph{held breath}, the fixed locus of the object under its own half-turn, which turns out to be exactly the torsion of the family.

\section{The forcing chain}

\begin{theorem}[the chain forces itself]
The following is a chain of forced determinations, each step canonical: \(\sigma \) (the substitution) forces \(\varphi \) (its Perron eigenvalue); \(\varphi \) forces \(A_{1} = \left(\begin{smallmatrix}2 & 1\\1 & 1\end{smallmatrix}\right)\) (the companion of \(x^{2} -\) 3x + 1, up to conjugacy) and its half \(X_{1} = \left(\begin{smallmatrix}1 & 1\\1 & 0\end{smallmatrix}\right); A_{1}\) forces \(4_{1}\) (its mapping torus); \(4_{1}\) forces the geometric representation \(\rho _{g}eo; \rho _{g}eo\) forces the invariant trace field \(\mathbb{Q}(\sqrt{-3})\). No step carries a choice. The object is not assembled; it is deduced from two arrows.
\end{theorem}

The metallic family is \(\sigma 's\) one-parameter deformation: \(\sigma _{m}\): a \(\to  a^{m}b\), b \(\to \) a, with eigenvalue the metallic mean \(\lambda _{m} = (m + \sqrt{m^{2}+4})/2\), monodromy \(A_{m} = R^{m} L^{m}\), and half \(X_{m} = \left(\begin{smallmatrix}m & 1\\1 & 0\end{smallmatrix}\right)\).

\section{The norm identity, and its freezing}

\begin{theorem}[the norm identity]
For every m, det \(X_{m} = N(\lambda _{m}) = -1\): the orientation character of the metallic half and the field norm of the metallic scale coincide, because \(X_{m}\) is the companion matrix of \(x^{2} -\) mx \(-\) 1, whose constant term is \(-1\). The half is orientation-reversing at every member; the scale is a unit of norm \(-1\) at every member.
\end{theorem}

\begin{theorem}[freezing through the golden degeneracy]
The members m \(= 1\), 4, 11, 29, … — the odd powers of \(\varphi  (\lambda _{m} = \varphi ^\{1,3,5,7,…\})\) — all live in the same field \(\mathbb{Q}(\sqrt{5})\), and the norm \(-1\) is carried unbroken through the degeneracy. No imaginary quadratic field admits a unit of norm \(-1\); the residue is therefore a real-field phenomenon, and it cannot be quantized away.
\end{theorem}

The norm identity is the residue in its arithmetic dress. It says the object's handedness is not an added structure but the \emph{norm of its own growth rate}, forced by the shape \(x^{2} -\) mx \(-\) 1 of the recurrence. You cannot have the golden growth without the negative norm; the sign is in the seed.

\section{The residue and the Pisano rhythm}

\begin{theorem}[the determinant law]
Define the residue of a word w \(\in \) \{R,L\}\emph{ as \(det(Par\cdot W(w))\) at level 15. Then residue(w) \(= -\omega ^\{\#L-\#R\}, \omega  = e^\{2\pi i/3\}\): it depends only on the letter imbalance mod 3, via the elementary determinants det \(W_{R} = \omega ̄\), det \(W_{L} = \omega \), det Par \(= -1\). Balanced words (\#R \(=\) \#L) have residue \(-1\) — }frozen*.
\end{theorem}

\begin{theorem}[the Pisano rhythm]
Along the Fibonacci letter tower (a \(\to \) R, b \(\to \) L), the imbalance at rung k is \(\pm F_\{k-2\}\), so the residue cycles with the Fibonacci sequence mod 3 — with period \(\pi (3) = 8\), the Pisano period. The residue does not merely survive the tower; it \emph{beats} along it, in the object's own arithmetic rhythm.
\end{theorem}

This is the residue in its quantum dress. The same \(\mathbb{Z}/2\) that is the norm of the scale is the determinant of the quantized word, and along the tower it acquires a heartbeat: a period-8 pulse set by the golden recurrence reduced mod the Markov prime 3.

\section{The two-register breath}

\begin{theorem}[the two registers agree]
The quantum orientation det(Par@N) and the classical orientation \(sign(\sigma \) acting on \((\mathbb{Z}/N)^{2}\)) are equal, and both equal \((-1)^\{(N-1)/2\}\) (Jacobi). At the seam levels 15, 45, 75, 225 they breathe together: the residue read through the quantized operator and the residue read through the classical map are the same sign, at every level.
\end{theorem}

The breath is the residue seen in two mirrors at once — the quantum and the classical — and the theorem is that the mirrors agree. There is one orientation, and the two quantizations of the object are two ways of reading it, never two values.

\section{The held breath is torsion}

The residue is what survives the object's motion. The final theorem asks what survives the object \emph{holding still} — the fixed locus under its own half-turn \(\sigma _{m}\).

\begin{theorem}[held breath \(= torsion\)]
On the cusp locus \(\kappa  = -2\), the characters fixed by \(\sigma _{m}\) are exactly the order-d torsion characters for divisors d \(\ge \) 3 of m. The geometric (hyperbolic) character is never fixed — it always breathes, exchanged by \(\sigma _{m}\). What can hold still is precisely the torsion: a generator of finite order d, with d \(| m\), whose finite order makes \(\sigma _{m}\) act as the swap a \(\leftrightarrow \) b, whose fixed locus is the symmetric characters. \textbf{Member m holds the breath of every divisor d \(\ge \) 3 of m.} The field of the held breath is \(\Delta _{d} = \tau _{d}^{2}(\tau _{d}^{2} -\) 8), \(\tau _{d} = 2cos(2\pi /d)\): 3 \(| m\) gives \(\mathbb{Q}(\sqrt{-7})\) (order 3), 5 \(| m\) gives \(\mathbb{Q}(\sqrt{41})\) (order 5); m \(\in \) \{1,2,4\} hold no breath at all.
\end{theorem}

The held breath is the object's orbifold skeleton — the elliptic points where a generator becomes a rotation. The geometry breathes; the torsion holds. And the divisibility is exact: what a member can hold still is the torsion of its divisors, no more and no less.

\section{Method, mirrored}

The program's method is its subject. To find the residue we did what \(\sigma \) does: we generated everything the object touches and let the magnitudes cancel — the traces, the fields, the volumes, the values — and watched for what would not go. What would not go was the sign. The residue is not a result we added at the end; it is the fixed point of the whole procedure, the \(\mathbb{Z}/2\) that survived the arithmetic (the norm), the quantization (the determinant), the two registers (the breath), and the object's own stillness (the held breath). Cancel everything; what remains is real.

\section{What is classical, and what remains open}

Classical: units of norm \(-1\) in real quadratic orders (Dirichlet); the Pisano period (Lagrange); the Jacobi-symbol parity of the metaplectic determinant; the fixed-point theory of mapping-class actions on character varieties. Ours (scoped, pending the gate): the three-register identification of the residue (norm \(= determinant = parity\) as one \(\mathbb{Z}/2\)); the Pisano-8 beat of the quantized word; the held-breath-is-torsion theorem with its closed-form fields and the m \(\in \) \{1,2,4\} exception. Open: (i) the standard-word lemma (chain words word-mirror for all k); (ii) the heredity of the chain's torsion primes; (iii) whether the held-breath fields \(\mathbb{Q}(\sqrt{-7}), \mathbb{Q}(\sqrt{41})\), … carry any further structure beyond the order-d mechanism (each appearance requires its own reason — none is banked as cumulative).

\section{Appendix: reproducibility}

Norm identity \texttt{br\_n\_norm.py} (B469); the determinant law and Pisano rhythm \texttt{rf3\_quantum.py} (B470); the two-register breath \texttt{br1\_br2.py} (B469); the held-breath theorem \texttt{B479\_held\_breath\_torsion/held\_breath\_capstone.py} (verified to order-13). Each lock re-run at draft. The banked \(\lambda _{c}hain\) value was independently re-verified as a genuine limit (not a finite-rung artifact) during this paper's assembly; recorded in the corrections ledger.

\begin{thebibliography}{99}
\bibitem{markov} A.~A. Markoff, \emph{Sur les formes quadratiques binaires ind\'efinies}, Math. Ann. \textbf{15} (1879), 381--406.
\bibitem{goldman} W.~M. Goldman, \emph{The modular group action on real SL(2)-characters of a one-holed torus}, Geom. Topol. \textbf{7} (2003), 443--486.
\bibitem{cl} S.~Cantat and F.~Loray, \emph{Dynamics on character varieties and Malgrange irreducibility of Painlev\'e VI}, Ann. Inst. Fourier \textbf{59} (2009), 2927--2978.
\bibitem{dirichlet} P.~G.~L. Dirichlet, \emph{Vorlesungen \"uber Zahlentheorie}, 1863.
\bibitem{hb} J.~H. Hannay and M.~V. Berry, \emph{Quantization of linear maps on a torus}, Physica D \textbf{1} (1980), 267--290.
\end{thebibliography}

\end{document}

